%-*-latex-*-
%-*-latex-*-
\newcommand\COURSE{ciss245}
\newcommand\ASSESSMENT{q6002}
\newcommand\ASSESSMENTTYPE{Quiz}
\newcommand\POINTS{\textwhite{xxx/xxx}}

\input{myquizpreamble}
\input{yliow}
\input{\COURSE}
%-*-latex-*-

%-*-latex-*-

%-*-latex-*-
\renewcommand\TITLE{\ASSESSMENTTYPE \ \ASSESSMENT}

\renewcommand\EMAIL{}
\renewcommand\AUTHOR{eagraham1@cougars.ccis.edu}

\textwidth=6in
\begin{document}
\topmatterthree

%------------------------------------------------------------------------------
\nextq
Write a class \verb!Time! that has private members
\verb!hh_!,
\verb!mm_!,
\verb!ss_!,
for hours, minutes, seconds in 24-hour format, i.e.,
hours is in the range 0..23.
Write enough public methods so that you can
create an array \verb!t! of 10 \verb!Time! objects
so that
\verb!t[0]! models the time 00:00:00,
\verb!t[1]! models the time 00:01:00,
\verb!t[2]! models the time 00:02:00,
...
\verb!t[9]! models the time 00:0\textred{9}:00.
Print all the object\textred{s} in \verb!t!.
\\
\ANSWER
\begin{answercode}
#include <iostream>

class Time
{
public:
void set_time(int hh, int mm, int ss)
{
  hh_ = hh;
  mm_ = mm;
  ss_ = ss;
}

void println() const
{
  std::cout << std::setw(2) << std::setfill('0') << hh_ << ":"
  << std::setw(2) << std::setfill('0') << mm_ << ":"
  << std::setw(2) << std::setfill('0') << ss_ << std::endl;
}
private:
int hh_;
int mm_;
int ss_;
};

int main()
{
    // Declare t to be an array of 10 Time objects
    Time t[10];
    // In a loop, let t[i] invoke method setHHMMSS(hh, mm, ss)
    // with the right values for hh, mm, ss.
    for (int i = 0; i < 10; i++)
    {
      time[i].set_time(0, i, 0);
    }
    // In a loop, let t[i] invoke method println() to print
    // the values in t[i] in the format hh:mm:ss, followed by newline.
    // Note that hh:mm:ss must be zero-filled. For instance
    // midnight must be printed by 00:00:00 and not 0:0:0.

    for (int i = 0; i < 10; ++i)
    {
      t[i].println();
    }
    
    return 0;
}
\end{answercode}

%------------------------------------------------------------------------------
\nextq
This is the same as the previous question \textit{except}
that \verb!t! is an array of \verb!Time! pointers.
You can copy the \verb!Time! class from above to the answer here.
\\
\ANSWER
\begin{answercode}
#include <iostream>

class Time
{
public:
void set_time(int hh, int mm, int ss)
{
  hh_ = hh;
  mm_ = mm;
  ss_ = ss;
}

void println() const
{
  std::cout << std::setw(2) << std::setfill('0') << hh_ << ":"
  << std::setw(2) << std::setfill('0') << mm_ << ":"
  << std::setw(2) << std::setfill('0') << ss_ << std::endl;
}
private:
int hh_;
int mm_;
int ss_;
};

int main()
{
    // Declare t to be an array of 10 Time pointers.
    // In a loop, allocate memory on the heap for t[i] to point to.
    Time * t[10]
    // In a loop, let the object that t[i] points to invoke method
    // setHHMMSS(hh, mm, ss) with the right values for hh, mm, ss.
    for (int i = 0; i < 10; i++)
    {
      time[i]->set_time(0, i, 0);
    }
    // In a loop, let the object t[i] points to invoke method println()
    // to print the values that the the object t[i] points.
    // Print in the format hh:mm:ss, followed by newline.
    // Note that hh:mm:ss must be zero-filled. For instance
    // midnight must be printed by 00:00:00 and not 0:0:0.
    for (int i = 0; i < 10; ++i)
    {
      t[i]->println();
    }
    // Deallocate memory used by t.
    for (int i = 0; i < 10; i++)
    {
      delete t[i];
    }
    
    return 0;
}

\end{answercode}

%------------------------------------------------------------------------------
\nextq
Wumpus lives in the Wumpus world, which is a 4-by-4 grid.
For instance in the following
\begin{console}[frame=single,fontsize=\footnotesize]
+-+-+-+-+
| | | | |
+-+-+-+-+
| | | |W|
+-+-+-+-+
| | | | |
+-+-+-+-+
| | | | |
+-+-+-+-+
\end{console}
Wumpus is at row 1, column 3.
You are given this:
\begin{console}[frame=single,fontsize=\footnotesize]
// file: WumpusWorld.h
#ifndef WUMPUSWORLD_H
#define WUMPUSWORLD_H

#include <iostream>

class WumpusWorld
{
public:
    void init();         // initialize so that Wumpus is at (0,0).
    void println();      // print according to the above format
    void move_wumpus();  // randomize a direction rand() % 4 where 0,1,2,3
                         // are N,S,E,W for Wumpus. Of course Wumpus
                         // must stay in the world.
                         // If wumpus tries to move N but its row is 0,
                         // wumpus stays put.
private:
    char world_[4][4];
};
#endif
\end{console}
Do not add anything to the above.
You want to make Wumpus wonder around like this:
\begin{console}[frame=single,fontsize=\footnotesize]
#include <iostream>
#include <ctime>
#include <cstdlib>
#include "WumpusWorld.h"

int main()
{
    srand((unsigned int) time(NULL));

    WumpusWorld ww;
    ww.init();
    ww.println();
    for (int i = 0; i < 5; ++i)
    {
        ww.move_wumpus();
        ww.println();
    }
    return 0;
}
\end{console}
Complete the following cpp file that contains the implementation
of the methods declared in the WumpusWorld class.
\\
\ANSWER
\begin{answercode}
// file: WumpusWorld.cpp

void WumpusWorld::init()
{
  for (int i = 0; i < 4; ++i)
  {
    for (int j = 0; j < 4; ++j)
    {
      world_[i][j] = ' ';
    }
  }

  world_[0][0] = 'W';
}

void WumpusWorld::println()
{
  for (int i = 0; i < 4; i++)
  {
    std::cout << "+-+-+-+-+" << std::endl;
    for (int j = 0; j < 4; j++)
    {
      std::cout << "|" << world_[i][j];
    }
    std::cout << "|" << std::endl;
  }
  std::cout << "+-+-+-+-+" << std::endl;
}

void WumpusWorld::move_wumpus()
{
  int wumpus_row = -1, wumpus_col = -1;
  for (int i = 0; i < 4; i++)
  {
    for (int j = 0; j < 4; j++)
    {
      if (world_[i][j] == 'W')
      {
        wumpus_row = i;
        wumpus_col = j;
        break;
      }
    }
    if (wumpus_row != -1)
    {
      break;
    }
  }

  int direction = rand() % 4; // 0 - North, 1 - South, 2 - East, 3 - West

  world_[wumpus_row][wumpus_col] = ' ';

  switch (direction)
  {
    case 0: // North
    if (wumpus_row > 0)
    {
      wumpus_row--;
    }
    break;
    case 1: // South
    if (wumpus_row < 3)
    {
      wumpus_row++;
    }
    break;
    case 2: // East
    if (wumpus_col < 3)
    {
      wumpus_col++;
    }
    break;
    case 3: // West
    if (wumpus_col > 0)
    {
      wumpus_col--;
    }
    break;
  }

  world_[wumpus_row][wumpus_col] = 'W';
}

\end{answercode}

%------------------------------------------------------------------------------
%-*-latex-*-
\newpage
\input{instructions}
\end{document}

